\section{附录：由对应构造的一个范畴}
设 $C\rightrightarrows X$ 是空间的群胚。
本附录（\S \ref{SS: cat of corr}）的主要目标是引入范畴 $\on{D}^C(X)$：当 $C=X\times_YX$ 且 $X\to Y$ 紧合时，它抓住了 $\on{D}(Y)$ 的大部分信息。

固定完全域 $k$，记 $\mathrm{pt}$ 为 $\Spec k$。
所谓叠（相应地，完全叠）$\mX$，指 $\Aff_k$（相应地，$\Aff_k^\pf$）上带 fpqc 拓扑的叠，其对角线可由代数空间（相应地，完全代数空间）表示，且存在从概形（相应地，完全概形）$X\to \mX$ 的平滑（相应地，完全平滑）满射。
后文中，$k$ 上（完全）有限表现的（完全）叠也称空间。关于完全代数几何基础的一些讨论，参看 \cite[附录 A]{Zh1}、\cite{BS} 与 \cite[\S A.1]{XZ}。

\subsection{上同调对应的回顾}
\label{Sec:cohomological correspondence}
我们大量用到上同调对应的形式化，故按 \cite[\S A.2]{XZ} 简要回顾若干事实。更详尽的处理见文献同前处。

设 $E$ 是系数环，$\cha k$ 在 $E$ 中可逆。
若 $X$ 是 $k$ 上（完全）有限表现概形且 $E$ 有限，记 $\on{D}(X)=\on{D}(X,E)$ 为 $E$-系数 étale 层（无界）导出范畴。对一般空间 $X$ 与可能更一般的系数环 $E$（例如 $E=\Ql$），可通过某种极限定义 $\on{D}(X)=\on{D}(X,E)$（精确表述见 \cite[\S A.1.15]{XZ}）。
记 $\on{D}^*_c(X)\subset \on{D}(X,E)$（$*=+,-,b,\emptyset$）为其可构造导出范畴，$\on{P}(X)\subset\on{D}_c^b(X)$ 为反常层范畴。记 $\omega_X\in \on{D}_c^b(X)$ 为 $X$ 上的对偶层，$\bD\mF=R\underline{\Hom}(\mF,\omega_X)$ 为 Verdier 对偶。按定义有典范赋值映射
\begin{equation}
\label{E: VD counit}
 \ev_\mF: \Delta^*(\bD\mF\boxtimes\mF)=\bD\mF\otimes \mF\to \omega_X,
\end{equation}
且若 $\mF\in\on{D}_c^b(X)$，取 Verdier 对偶即得典范余赋值映射
\begin{equation}
\label{E: VD unit}
\coev_\mF: E\to \Delta^!(\mF\boxtimes\bD\mF),\quad
\end{equation}

下面引入的概念很有用。
\begin{dfn}
\label{D: basechangeable}
空间交换方块
\begin{equation}
\label{E: basechangeable}
\begin{CD}
D@>p>>C\\
@VbVV@VVaV\\
Y@>f>>X
\end{CD}
\end{equation}
称为\emph{可基变换}\zy{[译者按（上下文推断，无外部依据）：base changeable 是原作者自定的术语，取「可基变换」，指该方块在任意基变换下仍给出同构。]}，若诱导态射 $h:D\to C\times_XY$ 可由（完全）紧合代数空间表示。
\end{dfn}

下面的引理是直接的。
\begin{lem}
\label{L: criterion basechangeable}
\begin{enumerate}
\item 如 \eqref{E: basechangeable} 的交换图可基变换，若
\begin{itemize}
\item $p$ 可由（完全）紧合代数空间表示且 $f$ 分离；
\item 或 $b$ 可由（完全）紧合代数空间表示且 $a$ 分离。
\end{itemize}
\item 在交换图
\begin{equation}
\label{E: two base changeable}
\begin{CD}
E@>q>>D@>p>>C\\
@VcVV@VbVV@VVaV\\
Z@>g>>Y@>f>>X
\end{CD}
\end{equation}
中，若两个内方块均可基变换，则外矩形也可基变换。
\end{enumerate}
\end{lem}

这一命名由下述事实支持。

\begin{lem}
\label{L:proper base change}
假设有如上 \eqref{E: basechangeable} 的可基变换空间交换图。
则存在自然的基变换同态
\begin{equation}
\label{E: base change isom}
\on{BC}^*_!: a^*f_!\to p_!b^*,
\end{equation}
若 \eqref{E: basechangeable} 是笛卡儿方块，则它是同构。
若 $p$ 与 $f$ 分离，它进入交换图
\[\begin{CD}
a^*f_!@>>> p_!b^*\\
@VVV@VVV\\
a^*f_*@>>>p_*b^*,
\end{CD}\]
其中下方箭头是自然伴随。
此外，给定 \eqref{E: two base changeable} 且两方块均可基变换，基变换同态 $a^*(fg)_!\to (pq)_!c^*$ 等于复合
$$a^*(fg)_!=a^*f_!g_!\xrightarrow{\on{BC}_!^*} f_!b^*g_!\xrightarrow{\on{BC}_!^*} f_!g_!c^*=(fg)_!c^*.$$
同样，$(fg)^*a_!\to c_!(pq)^*$ 等于复合
$$(fg)^*a_!=g^*f^*a_!\xrightarrow{\on{BC}_!^*} g^*b_!f^*\xrightarrow{\on{BC}_!^*} c_!g^*f^*=c_!(fg)^*.$$
\end{lem}

\begin{definition}
\label{AD:correspondences}
设 $(X_i,\mF_i)$（$i=1,2$）是两对数据，其中 $X_i$ 是空间，$\mF_i\in \on{D}(X_i,E)$。\emph{上同调对应} $(C,u):(X_1,\mF_1)\to (X_2,\mF_2)$ 指空间 $C\xrightarrow{c_1\times c_2} X_1\times  X_2$ 连同一个态射 $u:c_1^*\mF_1\to c_2^!\mF_2$。称 $C$ 为 $(C,u)$ 的\emph{支撑}。为简单起见，$(C,u)$ 有时也记作 $C$ 或 $u$。两个上同调对应 $(C,u)$ 与 $(C',u')$（从 $(X_1,\mF_1)$ 到 $(X_2,\mF_2)$）之间的（同）构有显然的概念。从 $(X_1,\mF_1)$ 到 $(X_2,\mF_2)$、支撑在 $C$ 上的上同调对应的同构类之集记作 $\on{Corr}_C((X_1,\mF_1),(X_2,\mF_2))$，按定义
\[\on{Corr}_C((X_1,\mF_1),(X_2,\mF_2))\cong \Hom_C(c_1^*\mF_1,c_2^!\mF_2).\]

若 $(C,u):(X_1,\mF_1)\to (X_2,\mF_2)$ 与 $(D,v):(X_2,\mF_2)\to (X_3,\mF_3)$ 是两个上同调对应，定义其复合为 $(C\times_{X_2}D,v\circ u):(X_1,\mF_1)\to (X_3,\mF_3)$，其中 $v\circ u$ 是如下复合
\[p^*c_1^*\mF_1\xrightarrow{p^*u} p^*c_2^!\mF_2\xrightarrow{\on{BC}^{*!}} q^!d_1^*\mF_2\xrightarrow{q^!v} q^!d_2^!\mF_3,\]
这里 $\on{BC}^{*!}$ 是 \eqref{E: base change isom} 中的基变换同态，$p,q$ 分别是 $C\times_{X_2}D$ 到 $C$ 与 $D$ 的投影。
\end{definition}

下面是一些例子。

\begin{ex}
\label{Ex:examples of correspondences}
\begin{enumerate}
\item 对态射 $f: X\to Y$ 与 $\mF \in \on{D}(X)$，有上同调对应 $(X\xleftarrow{\id}X\xrightarrow{f} Y, u:\mF\to f^!f_!\mF)$，从 $(X,\mF)$ 到 $(Y,f_!\mF)$，称为\emph{前推对应}，简记作 $(\Ga_f)_!$。

\item 对态射 $f: X\to Y$ 与 $\mF \in \on{D}(Y)$，有自然的上同调对应 $(Y\xleftarrow{f} X\xrightarrow{\id} X, u:= \mathrm{id}:f^*\mF\to f^*\mF)$，从 $(Y,\mF)$ 到 $(X,f^*\mF)$，称为\emph{拉回对应}，简记作 $\Ga_f^*$。

\item 若 $(C,u): (X_1,\mF_1)\to (X_2,\mF_2)$ 是上同调对应，且 $\mF_i\in \on{D}_c^b(X_i)$，则 $(C,\bD u): (X_2, \bD \mF_2)\to (X_1,\bD \mF_1)$ 也是上同调对应，称为\emph{对偶对应}。

\item 设 $(X,\mF)$ 是一对数据，$\mF\in \on{D}_c^b(X)$。则 \eqref{E: VD counit} 与 \eqref{E: VD unit} 给出上同调对应及其对偶
\[\begin{aligned}
&\ev_\mF\in \on{Corr}_{X}((X\times X, \bD\mF\boxtimes \mF),(\on{pt}, E)),\\
&\coev_\mF\in \on{Corr}_{X}((\on{pt}, E), (X\times X, \mF\boxtimes\bD\mF)).
\end{aligned}\]

\item 假设 $k=\bar k$。设 $X_1\leftarrow C\to X_2$ 是对应。则
\begin{equation*}
\label{E:corr and cycle}
\begin{split}
\on{Corr}_C((X_1,E[d_1]), (X_2, \omega_{X_2}[d_2]))&=\Hom_{\on{D}^b_c(C)}(E[d_1],\omega_C[d_2])\\
&=\on{H}^{\on{BM}}_{d_1-d_2}(C).
\end{split}
\end{equation*}
故若 $2\dim C=d_1-d_2$，则 $\on{Corr}_C((X_1,E[d_1]), (X_2, \omega_{X_2}[d_2]))$ 可等同于 $C$ 的最高次 Borel–Moore 同调。特别地，若 $X_1,X_2,C$ 是有限集，看作 $k$ 上的零维空间，则 $\on{Corr}_C((X_1,E),(X_2,E))$ 可等同于 $C$ 上 $E$ 值函数之空间。
\end{enumerate}
\end{ex}

现假设有交换图
\[\begin{CD}
X_1@<c_1<< C@>c_2>> X_2\\
@Vf_1VV@VVfV@VVf_2V\\
Y_1@<d_1<<D@>d_2>>Y_2
\end{CD}\]
且设 $(C, u:c_1^*\mF_1\to c_2^!\mF_2)$ 是从 $(X_1,\mF_1)$ 到 $(X_2,\mF_2)$ 的上同调对应。
设左交换方块可基变换（定义 \ref{D: basechangeable}），则有
\begin{equation}
\label{ASS:pushforward correspondence}
d_1^*(f_1)_!\mF_1\xrightarrow{\on{BC}^*_!} f_!c_1^*\mF_1\xrightarrow{f_!u} f_!c_2^!\mF_2\to d_2^!(f_2)_!\mF_2,
\end{equation}
其中最后一个映射是自然伴随。
为记号简便，上述映射仍记作 $f_!(u)$。则 $(D,f_!(u))$ 是从 $(Y_1,(f_1)_!\mF_1)$ 到 $(Y_2,(f_2)_!\mF_2)$ 的上同调对应。称它为\emph{上同调对应的前推}。

特别地，若 $(Y_1\xleftarrow{d_1} D\xrightarrow{d_2} Y_2)$ 为 $(S=S=S)$（$S$ 是某基概形），且 $c_1$ 紧合（故左方块可基变换），则得到层的同态
\[f_!(u): (f_1)_!\mF_1\to (f_2)_!\mF_2.\]
更具体地，若 $S=\Spec k$ 代数闭且 $c_1$ 紧合，则得到上同调上的诱导映射 $\on{H}_c(u): \on{H}_c^*(X_1,\mF_1)\to \on{H}_c^*(X_2,\mF_2)$。

\begin{rmk}
若 $f_1$ 分离且 $f$ 可由完全紧合代数空间表示（故左方块可基变换），则 $f_!(u)$ 分解为自然映射 $(f_1)_!\mF_1\to (f_1)_*\mF_1$ 与
\[d_1^*(f_1)_*\mF_1\to f_*c_1^*\mF_1\to f_!c_2^!\mF_2\to d_2^!(f_2)_!\mF_2\]
的复合。后一映射即 SGA 5 中考虑的上同调对应的前推。不过在我们处理的大多数情形中，$d_1$ 分离而 $c_1$ 可由完全紧合代数空间表示。
\end{rmk}
